\begin{abstract}
An agent that optimizes compute kernels faces an objective whose two halves pull
apart: many of the cheapest ways to make a kernel faster also make it wrong. The
incorrect ``optimizations'' we timed run up to \sabSpeedNoDivide$\times$ faster
than the correct kernel, and a constant-fill ``softmax'' runs \fillSpeedup$\times$
faster \emph{and passes our own property gate}. The correctness gate is therefore
half of the fitness function, and its calibration and completeness bound what a
long-running agent can reach. We describe a harness with three layers: agent
\emph{skills} that carry measured lessons across sessions and control what each
sub-agent may see; \emph{property-based testing} over byte-identical replayed
executions, which serves as a cheap and calibrated gate; and \emph{formal proof},
which certifies that the gate's laws pin down the specification. Comparing four oracle
strategies on NumPy and Triton (\gpuName), we find that they never disagree on
\mutantsAdmitted{} mutants from three corpora; that the field-default \allclose
rejects \fpAllcloseRate{} of correct-kernel groups while the others reject none,
so an \allclose-gated loop cannot accept its own baseline; and that a
property-gated loop improves validation bits-per-byte by \loopBpbGain\% at
\loopTputGain\% higher throughput. We also show that the property bundles are
incomplete: they admit \gapFamiliesTotal{} wrong kernels, which we characterize
exactly and close with one law per task. Last, we describe the formal layer. An
agent writes the spec, a model of the kernel and the proof in each of seven formal
languages, and six checks decide whether the proof certifies anything. Those checks
execute every formal artifact on the PBT cases and catch each kind of spec gaming.
Their core is built, and the pilot is pre-registered.
\end{abstract}
